\chapter{Um Primeiro Mergulho em JavaScript}

\section{O projeto: jogo de adivinhação de números}

A ideia do jogo é simples de descrever, mas rica o suficiente para ilustrar praticamente todos os conceitos fundamentais da linguagem. Quando a página carrega, o programa sorteia um número aleatório entre 1 e 100. O usuário tem 10 tentativas para descobrir esse número, digitando um valor em um campo de texto e confirmando com um botão. A cada tentativa, o sistema informa se o palpite foi maior ou menor que o número sorteado. Se o usuário acertar dentro das 10 tentativas, uma mensagem de sucesso é exibida; caso contrário, o jogo termina e é exibido um botão para reiniciar.

\begin{infobox}{Requisitos funcionais do jogo}
• Sortear um número aleatório entre 1 e 100 no início do jogo.
\end{infobox}

\begin{itemize}
  \item Permitir até 10 tentativas de adivinhação.
\end{itemize}

\begin{itemize}
  \item A cada tentativa, informar se o valor digitado foi muito alto ou muito baixo.
\end{itemize}

\begin{itemize}
  \item Exibir o histórico de tentativas anteriores.
\end{itemize}

\begin{itemize}
  \item Ao acertar, exibir mensagem de parabéns e encerrar o jogo.
\end{itemize}

\begin{itemize}
  \item Ao esgotar as tentativas sem acertar, exibir mensagem de fim de jogo.
\end{itemize}

\begin{itemize}
  \item Oferecer um botão para reiniciar o jogo a qualquer momento após o término.
\end{itemize}

Esse exemplo é interessante justamente por ser pequeno o bastante para caber neste capítulo, mas completo o bastante para exercitar variáveis, constantes, funções, eventos, laços, manipulação do DOM e um pouco de orientação a objetos. Cada um desses conceitos será aprofundado nos capítulos seguintes; aqui, o objetivo é apenas ganhar familiaridade prática com a ”sensação”de programar em JavaScript.

\section{Estrutura HTML de partida}

Antes de escrever qualquer JavaScript, é preciso ter uma estrutura HTML mínima para trabalhar. O esqueleto do jogo é formado por um título, um parágrafo de instruções, um formulário simples (com um campo de texto e um botão) e uma área destinada às mensagens de status:

\begin{codebox}{HTML}
<!DOCTYPE html>
<html lang="pt-br">
<head>
  <meta charset="UTF-8">
  <title>Jogo de Adivinhação</title>
</head>
<body>
  <h1>Jogo de Adivinhação de Números</h1>
  <p>Selecionamos um número aleatório entre 1 e 100.
     Você consegue adivinhar em 10 tentativas ou menos?</p>

     <form>
       <label for="guessField">Digite uma tentativa: </label>
       <input type="text" id="guessField" class="guessField">
       <input type="submit" value="Enviar tentativa" class="guessSubmit">
     </form>

     <div class="resultParas">
       <p class="guesses"></p>
       <p class="lastResult"></p>
       <p class="lowHigh"></p>
     </div>

  <script src="jogo.js" defer></script>
</body>
</html>
\end{codebox}

Note que os elementos que precisarão ser manipulados por JavaScript recebem uma classe (ou poderiam receber um id) para que seja possível localizá-los depois a partir do código.

\section{Declarando as variáveis do jogo}

O primeiro passo dentro do arquivo JavaScript é declarar as variáveis que o programa vai utilizar. Uma variável é, essencialmente, um espaço reservado na memória do computador para armazenar um valor.

\begin{codebox}{JavaScript}
 let randomNumber = Math.floor(Math.random() * 100) + 1;
 let guessCount = 1;
 let resetButton;

 const guesses = document.querySelector(".guesses");
 const lastResult = document.querySelector(".lastResult");
 const lowOrHigh = document.querySelector(".lowHigh");

 const guessSubmit = document.querySelector(".guessSubmit");
 const guessField = document.querySelector(".guessField");

 guessField.focus();
\end{codebox}

Vale destacar o funcionamento de Math.random(), um método da classe Math que gera um número decimal aleatório entre 0 (inclusive) e 1 (exclusive). Multiplicando esse valor por 100, obtém-se um número entre 0 e 99,999...; aplicando Math.floor(), o valor é arredondado para baixo, resultando em um número inteiro entre 0 e 99. Somando 1, finalmente, chega-se a um número inteiro entre 1 e 100, exatamente como desejado. As constantes guesses, lastResult e lowOrHigh são ponteiros (referências) para os três parágrafos definidos dentro da div de resultado. Da mesma forma, guessSubmit e guessField apontam, respectivamente, para o botão e para o campo de texto do formulário. Uma vez obtidas essas referências, é possível manipular esses elementos livremente a partir do código.

\begin{infobox}{Dica}
Repare que querySelector recebe como argumento um seletor CSS — os mesmos seletores usados nas folhas de estilo. Um ponto na frente do nome (.guesses) indica que se está buscando um elemento pelo valor do atributo class; já um sustenido (\#jogador) indica busca por id. Essa é uma das grandes vantagens de já conhecer CSS antes de aprender JavaScript.
\end{infobox}

\section{Escrevendo a primeira função}

Uma função em JavaScript agrupa um bloco de código que pode ser executado (chamado) quando necessário. Para um primeiro teste, é possível escrever uma função simples que apenas exibe um alerta:

\begin{codebox}{JavaScript}
 function checkGuess() {
   alert("Eu sou um placeholder!");
 }
\end{codebox}

A palavra reservada function inicia a definição, seguida do nome escolhido para a função, parênteses (que conteriam os parâmetros, se houvesse algum) e um bloco de código delimitado por chaves. A função alert() é uma função global do navegador

que exibe uma caixa de diálogo simples com uma mensagem informativa — útil para testes rápidos, embora pouco usada em interfaces reais e profissionais.

\begin{infobox}{Dica}
Uma forma prática de testar pequenos trechos de JavaScript, sem precisar editar arquivos, é abrir o console do navegador (parte das ferramentas de desenvolvedor, acessíveis geralmente pela tecla F12 ou pelo menu de inspeção). É possível digitar expressões diretamente ali e ver o resultado na hora — inclusive chamar funções já definidas no script carregado pela página.
\end{infobox}

\section{Operadores básicos}

Antes de completar a lógica do jogo, vale revisar rapidamente os operadores mais comuns, todos testáveis diretamente no console do navegador:

\begin{codebox}{JavaScript}
6 + 9;       // 15   (adição)
20 - 15;     // 5    (subtração)
3 * 7;       // 21   (multiplicação)
10 / 5;      // 2    (divisão)

// concatenação de strings com o operador +
let name = "Cris";
"Olá, " + name;    // "Olá, Cris"

// operadores de comparação
5 === 2 + 3;   // true (estritamente igual: mesmo valor e mesmo tipo)
"2" === 2;     // false (tipos diferentes: string x number)
"2" !== 2;     // true (estritamente diferente)
\end{codebox}

Os operadores de comparação === (estritamente igual) e !== (estritamente diferente) levam em conta tanto o valor quanto o tipo dos operandos, e serão discutidos com mais detalhes no capítulo sobre números e operadores.

\section{Completando a lógica de verificação}

Com as variáveis e o entendimento básico de funções e operadores em mãos, é possível escrever a função que de fato verifica a tentativa do usuário:

\begin{codebox}{JavaScript}
function checkGuess() {
  const guess = Number(guessField.value);

     if (guessCount === 1) {
       guesses.textContent = "Tentativas anteriores: ";
     }



     guesses.textContent += guess + " ";

     if (guess === randomNumber) {
       lastResult.textContent = "Parabéns! Você acertou o número!";
       lastResult.style.backgroundColor = "green";
       lowOrHigh.textContent = "";
       setGameOver();
     } else if (guessCount === 10) {
       lastResult.textContent = "!!! GAME OVER !!!";
       setGameOver();
     } else {
       lastResult.textContent = "Errado!";
       lastResult.style.backgroundColor = "red";

         if (guess < randomNumber) {
           lowOrHigh.textContent = "Tentativa muito baixa!";
         } else if (guess > randomNumber) {
           lowOrHigh.textContent = "Tentativa muito alta!";
         }
     }

     guessCount++;
     guessField.value = "";
     guessField.focus();
 }
\end{codebox}

Repare em vários detalhes importantes deste trecho. Primeiro, guessField.value retorna o valor digitado pelo usuário no campo de texto — sempre como uma string — e por isso ele é convertido explicitamente para número com o construtor Number(). Segundo, a propriedade textContent permite ler e escrever o conteúdo textual de um elemento; já a propriedade style dá acesso às propriedades CSS daquele elemento diretamente pelo JavaScript, permitindo, por exemplo, mudar a cor de fundo dinamicamente. Por fim, a lógica condicional compara o palpite do usuário com o número sorteado e decide qual mensagem exibir.

\section{Capturando o evento de clique}

Até aqui, a função checkGuess existe, mas nada a executa automaticamente: é preciso registrar um ”ouvinte”do evento de clique no botão de envio:

\begin{codebox}{JavaScript}
 guessSubmit.addEventListener("click", checkGuess);
\end{codebox}

O método addEventListener recebe dois argumentos: o tipo do evento de interesse (aqui, "click", o clique do mouse) e a função que deve ser chamada quando o evento ocorrer. Note que o nome da função é passado sem parênteses — apenas uma referência à função, e não uma chamada imediata dela.

\section{Finalizando o jogo}

Falta ainda escrever a função que trata o fim do jogo, seja por vitória, seja por esgotamento das tentativas:

\begin{codebox}{JavaScript}
 function setGameOver() {
   guessField.disabled = true;
   guessSubmit.disabled = true;

     resetButton = document.createElement("button");
     resetButton.textContent = "Iniciar novo jogo";
     document.body.appendChild(resetButton);

     resetButton.addEventListener("click", resetGame);
 }
\end{codebox}

Esta função desabilita o campo de texto e o botão de envio (para impedir novas tentativas) e cria dinamicamente, via document.createElement, um novo botão de reinício, que é inserido no documento com appendChild. Em seguida, um evento de clique é associado a esse novo botão, apontando para uma função resetGame, responsável por reiniciar o estado do jogo:

\begin{codebox}{JavaScript}
 function resetGame() {
   guessCount = 1;

     const resetParas = document.querySelectorAll(".resultParas p");
     for (let i = 0; i < resetParas.length; i++) {
       resetParas[i].textContent = "";
     }

     resetButton.parentNode.removeChild(resetButton);

     guessField.disabled = false;
     guessSubmit.disabled = false;
     guessField.value = "";
     guessField.focus();

     lastResult.style.backgroundColor = "white";

     randomNumber = Math.floor(Math.random() * 100) + 1;
 }
\end{codebox}

Observe o laço for, que percorre todos os parágrafos de resultado limpando seu conteúdo, e a remoção do botão de reinício por meio de parentNode.removeChild() — já que um elemento não pode remover a si mesmo do DOM, mas seu elemento pai pode.

Síntese do Capítulo

\begin{itemize}
  \item O jogo de adivinhação de números serve como exemplo prático que reúne variáveis, funções, eventos, laços e manipulação do DOM em um único programa coeso.
\end{itemize}

\begin{itemize}
  \item Math.random() combinado com Math.floor() é a forma padrão de gerar números inteiros aleatórios dentro de um intervalo em JavaScript.
\end{itemize}

\begin{itemize}
  \item querySelector e querySelectorAll localizam elementos do DOM usando a mesma sintaxe de seletores do CSS.
\end{itemize}

\begin{itemize}
  \item textContent lê e altera o texto de um elemento; a propriedade style permite alterar CSS diretamente via JavaScript.
\end{itemize}

\begin{itemize}
  \item addEventListener associa uma função a um evento de um elemento, sendo a forma recomendada de tratar eventos (em vez de atributos como onclick).
\end{itemize}

\begin{itemize}
  \item Elementos podem ser criados dinamicamente com createElement e inseridos no DOM com appendChild, assim como removidos com removeChild a partir do elemento pai.
\end{itemize}

\begin{itemize}
  \item O console do navegador é uma ferramenta valiosa para testar pequenos trechos de código e funções isoladamente, antes de integrá-los ao programa completo.
\end{itemize}
